\section{THEORETICAL BACKGROUND}

\subsection{Historic Prospects Of TLP Tendon Modeling}

TLPs are floating offshore structures anchored to the seabed by vertical tendons maintained under high pre-tension. Their dynamic behavior results from the coupled interaction between hydrodynamic loads, tendon elasticity, and platform motion. Understanding these vibration characteristics is essential for ensuring structural integrity. Over the past decades, several studies have employed beam-based idealizations to represent the structural behavior of TLP components, aiming to simplify the analysis while retaining the essential dynamic characteristics of the system. Among the pioneering works, \textcite{Adrezin1999} developed an analytical model describing the stochastic response of a TLP subjected to irregular sea states.

In their formulation, the platform was treated as a rigid cylinder connected to a single elastic tendon represented by an Euler-Bernoulli beam. The study demonstrated that neglecting the tendon's flexural stiffness can significantly underestimate the overall dynamic response of the system, particularly in the surge motion. Their results highlighted the importance of including tendon dynamics in coupled platform analyses to obtain realistic estimates of vibration amplitudes and natural frequencies.

A broader theoretical overview is presented by \textcite{Chakrabarti2005}, where the author systematically discussed the assumptions and applicability of beam theory for representing TLP components. The author emphasized that while Euler-Bernoulli beam models provide valuable insight for linear vibration regimes, more complex behaviors, such as coupling between pitch and surge or hydrodynamic nonlinearity, require hybrid formulations combining beam and rigid-body dynamics. This distinction is critical for accurately predicting the global response of the platform when subjected to extreme environmental loading conditions.

Building on that premise, \textcite{oyejobi2018} investigated the dynamic coupling between the platform and its tendons through numerical simulations based on FEM. The tendons were modeled as tensioned beam elements, capturing geometric stiffness and hydrodynamic damping effects. The results confirmed that tendon flexibility and tension variations play a critical role in the system's dynamic amplification and in the distribution of stresses along the anchoring members. This finding reinforced the relevance of beam-based models as reliable tools for preliminary design and vibration prediction in TLPs.

The study by \textcite{Jahangiri2016}, proposed a diagnostic approach grounded in the vibration analysis of a piezoelectric energy harvesting system to identify structural damages, specifically reductions in stiffness and Young's modulus. Their results showed that even small deviations in modal properties could be effectively detected using fuzzy classification, reinforcing the practical importance of accurate dynamic modeling using beam analogies for TLP monitoring.

\textcite{MirandaJr2017} investigated the forced response and band structure of a 1D phononic crystal beam composed of steel and polyethylene, modeled using the EBT. The study employed a comprehensive comparative analysis of FEM, Spectral Element, and Wave-based formulations, alongside Plane Wave Expansion methods. The results demonstrated that the Improved Plane Wave Expansion method yields superior convergence compared to conventional approaches and experimentally validated the existence of low-frequency Bragg-type band gaps, indicating the potential of these periodic structures for effective vibration management.

More recent research has explored the applicability of EBT beams to the modeling of TLPs. \textcite{RamosJunior2022} performed a comparative study on the free vibration of TLP tendons using Euler-Bernoulli, Rayleigh, and Timoshenko beam theories implemented through the Finite Element Method. Their analysis revealed that for slender tendons under high tension, the Euler-Bernoulli model provides sufficiently accurate results, while for shorter or stiffer tendons, rotary inertia and shear deformation effects become relevant, and the Rayleigh and Timoshenko formulations yield improved accuracy. This comparison is particularly relevant to the present research, as it validates the adoption of the EBT model as an efficient and physically consistent representation for tendon-like elements in TLP systems.

Together, these studies provide a consistent foundation for the modeling TLP dynamics, using the EBT formulation. They demonstrate that such models capture the essential features of tendon vibration and coupling effects with satisfactory accuracy, serving as a robust framework for both analytical and numerical investigations. This body of research underpins the present study's choice to employ periodic Euler-Bernoulli beams as analog models to explore vibration attenuation mechanisms and band-gap formation in offshore structures.

%========================================================
\subsection{Application Of Metamaterials For Vibration Mitigation In TLPs}
%=======================================================
% transição conceitual entre TLP e vigas periódicas

The study of wave propagation in slender and tensioned structural members provides an intermediate theoretical foundation linking global platform dynamics to localized vibration phenomena. In tendon-like components, longitudinal and flexural waves can propagate along the structure, and their reflection, transmission, or attenuation depend strongly on geometric discontinuities and boundary conditions \cite{rao2007,meirovitch2001}.

In the context of offshore systems, the dynamic behavior of cables, risers, and tendons is frequently analyzed using tensioned beam models, which emphasize the critical influence of geometric stiffness on the structural response \cite{Adrezin1999}. Expanding on the mechanics of wave propagation in such beam-like structures, \textcite{Maurin2016} demonstrated how periodicity and deformation levels fundamentally alter wave dispersion and propagation speeds. These principles are particularly relevant to complex marine structures, where local variations in material or cross-sectional properties can act as filtering interfaces for vibration energy, a concept reinforced by the extensive work of \textcite{Senjanovic2010} on ship and offshore structural dynamics.

These findings have motivated the development of passive vibration control strategies based on engineered periodicity. Local modifications, such as stiffness variations, mass inclusions, or geometric transitions, can significantly alter the propagation characteristics of mechanical waves, leading to attenuation bands where excessive vibration is effectively suppressed. This concept, derived from the physics of phononic and elastic metamaterials, has been successfully applied in structural engineering to mitigate unwanted vibrations through geometry-induced band gaps \cite{hussein2014,MirandaJr2017}. 

Consequently, the application of periodicity-based principles to tendon-like elements offers a promising path for vibration mitigation in offshore platforms. The modeling of these systems through beam theories allows the identification of frequency ranges where flexural or longitudinal waves are attenuated, establishing a direct connection between the global dynamic behavior of TLPs and the localized wave phenomena explored in periodic structures.

Given the sensitivity of these slender components to dynamic excitation, standard structural damping often proves insufficient. To address this, recent research has turned to the concept of periodic structures. A periodic structure is defined as a physical system exhibiting repeating spatial patterns, which may manifest in its constituent materials, internal geometry, or boundary conditions. Historically, structural engineering has explored this concept in systems such as aircraft fuselages, ship hulls, and heat exchanger tube banks, where periodic structure theory simplified the analysis of dynamic response and reduced computational costs. In these classical applications, it was observed that such structures behave as band-pass filters, efficiently radiating noise or vibrating at specific frequencies while attenuating others. The fundamental characteristic of these systems is their unique wave propagation behavior, which can be analyzed by examining a single repeating unit cell \cite{hussein2014,SenGupta1980}.

These systems present the distinctive ability to inhibit wave propagation within specific frequency ranges, known as band gaps. In conventional periodic structures, this phenomenon is primarily driven by Bragg scattering, a mechanism where the spatial periodicity of material properties or geometry induces destructive interference of waves when the wavelength is comparable to the unit cell size . Building on these principles, researchers have adapted concepts from solid-state physics to design elastic metamaterials and employ ``substitute continuum'' approaches to simplify analysis. By replacing the discrete periodic structure with a continuum beam of equivalent properties, engineers can efficiently predict global dynamic behaviors, such as vibration attenuation governed by Bragg effects, using established mechanical theories \cite{MirandaJr2017,hussein2014}.

% The accuracy of these continuum analogies depends heavily on the specific beam theory applied. In the analysis of one-dimensional periodic trusses, \textcite{Stephen1999} demonstrated that while the Euler-Bernoulli theory provides acceptable accuracy for the lowest modes, it significantly overestimates frequencies for higher modes where shear deformation becomes non-negligible. This limitation is further corroborated by \textcite{Domagalski2021}, who compared natural frequencies in periodic stepped and bi-material beams. His findings confirm that while Euler-Bernoulli assumptions hold for slender elements at low frequencies, they lead to unacceptable discrepancies in the higher frequency range compared to the Timoshenko-Ehrenfest theory. Since the formation of band gaps often involves shorter wavelengths where shear deformation and rotational inertia are critical, these studies underscore the necessity of verifying the validity limits of classical beam theories for accurate periodic design.

% The Euler-Bernoulli theory remains a fundamental model for studying wave propagation and vibration phenomena in slender structures. In the specific context of offshore engineering, \textcite{RamosJunior2022} validated this formulation for the free vibration analysis of Tension Leg Platform (TLP) tendons, demonstrating that for slender geometries under high tension, it yields results comparable to higher-order theories. However, when periodicity is introduced, either through geometry (stepped beams) or material distribution (bimaterial beams), the application of this theory requires careful consideration of its frequency limits.

% As demonstrated by \textcite{Domagalski2021}, while the Euler-Bernoulli model is highly accurate for slender periodic beams at lower frequencies, it diverges from the Timoshenko-Ehrenfest theory in the high-frequency regime where shear deformation and rotational inertia become significant. Notwithstanding, for the extremely slender geometry of TLP tendons, the Euler-Bernoulli formulation remains a robust and computationally efficient framework for capturing the relevant global structural modes and primary attenuation bands, offering a balance between physical accuracy and model complexity.

To characterize the vibration suppression capabilities of realistic structures, direct numerical approaches provide a crucial alternative to classical infinite-medium theories. As employed by \textcite{Domagalski2021} in the study of periodic beams, the numerical solution of the generalized eigenvalue problem via FEM constitutes a practical framework for identifying attenuation zones. This approach allows for the detection of band gaps by analyzing the distribution and clustering of natural frequencies, where stopbands manifest as spectral ranges devoid of resonance modes. Distinguishing itself from infinite-medium assumptions, this direct analysis enables the characterization of finite systems, facilitating a straightforward comparison with global structural requirements.

Hybrid approaches are required to investigate the compelx dynamic behavior of these systems. As detailed by \textcite{Reis2024} in the analysis of periodic structures, the Transfer Matrix Method combined with Bloch-Floquet theory constitutes a robust analytical framework for calculating the complex band structure of infinite periodic beams. This approach allows for the precise identification of attenuation zones (band gaps) based on the unit cell properties. Complementing this analytical insight, FEM is employed to simulate finite periodic beams, enabling the verification of mode shapes and the direct comparison of transmission spectra with practical structures. Together, these approaches form the theoretical backbone for studying periodic beams as analogs of vibration control systems.

% However, the application of the Euler-Bernoulli theory to periodic beams requires careful consideration of its frequency limits. As also discussed by \textcite{Domagalski2021} in the comparison of beam theories for periodic stepped structures, the Euler-Bernoulli model remains highly accurate for slender elements at lower frequencies but diverges from the Timoshenko-Ehrenfest theory in the high-frequency regime where shear deformation becomes significant. Notwithstanding, for the extremely slender geometry of TLP tendons ($L/d\gg 100$), the Euler-Bernoulli formulation remains a robust framework for capturing the relevant global structural modes and primary attenuation bands , balancing physical accuracy with the computational efficiency required for parametric design.

The transition from theoretical periodic beams to practical offshore applications is well-supported in the literature. \textcite{Asiri2014}, for instance, demonstrated that introducing periodicity through passive inserts in offshore pipes effectively attenuates wave propagation, significantly reducing vibration transmission. The efficiency of these attenuation bands depends strongly on the geometric parameters of the unit cell. As shown by \textcite{Iqbal2023} in parametric studies of periodic tracks, even subtle modifications in mass and stiffness distribution can lead to substantial changes in the attenuation characteristics. Within this context, the present research introduces a geometric saturation parameter, $\alpha$, to systematically investigate these sensitivities, linking the unit cell geometry directly to the band gap formation mechanisms essential for passive control in offshore structures.

% The efficiency of these attenuation bands depends strongly on the geometric and material parameters of the unit cell. Variables such as stiffness ratio, density contrast, segment length ratio, and periodicity spacing govern the onset and extent of forbidden frequency ranges. Recent parametric studies, such as those by \textcite{Iqbal2023} on supported tracks, reinforce that even subtle modifications in these properties, specifically the stiffness and mass distribution, can lead to substantial changes in the attenuation characteristics.

% Within this context, the present research introduces a geometric saturation parameter, denoted by $\alpha$, as a compact and systematic means of investigating how variations in unit cell configuration affect the vibrational response. This parameterized approach enables a clearer understanding of design sensitivities, linking the geometric definitions of the unit cell directly to the band gap formation mechanisms essential for optimized passive control in offshore structures.

Within this context, the present research evaluates the formation and location of band gaps via a saturation parameter $\alpha$. This parameter changes unit cell configuration to allow for modified structural periodicity, and subsequent tuning of band gap magnitude and distribution along the frequency spectrum. 

% Section 3 presents the EBT model developed via analytical and FEM formulations. While section 4 presents validations of relevant routines and simulation of a tendon-like structure. Finally, section 5 displays the main conclusions drawn from this research.